\begin{textsample}{NEG Dim 1 – vem\_ed\_09 – Score -1.00 – vem\_ed\_09/t038.txt}  \label{ex:f1_neg_014}
BOAS PRÁTICAS Competências culturais: além da disciplina Mestre em Economia Aplicada pela Esalq / USP, onde se graduou em Engenharia Agronômica, Ricardo Bertoni Pompeu é professor da Fatec Americana desde 2010 e atua com COIL / PCIs há quatro anos.
Pompeu compartilha um pouco de sua experiência no relato a seguir: Em 2017, participei de um megaprojeto com Eva Haug e Ariane Hoekstra, da Amsterdam University of Applied Sciences (AUAS / Holanda) e Osvaldo Succi Junior (CPS), Cristine Moraes (Fatec Piracicaba), Ana Teresa Colenci Trevelin (Fatec São Carlos), Carlos Amaral Moreira e eu (Fatec Americana).
Durante 7 semanas, 70 alunos estudaram estratégias de comunicação de marcas de suco de frutas, cerveja, refrigerante, chocolate, bolacha e sorvete nos dois países.
Em 2018, nosso projeto envolveu os professores Arun Pillutla (St.
Ambrose University / EUA), Carlos Moreira, Osvaldo Succi Junior e eu.
O PCI abordou a competência de dar e receber feedback, com a Metodologia Ativa de Aprendizado Baseado em Problemas.
Os alunos receberam uma situação problema e textos teóricos sobre feedback, incluindo o aspecto cultural dessa competência.
Em 2019 tivemos um novo PCI com a Holanda, com Eva Haug e Brechtine Detmar (AUAS) e Osvaldo Succi Junior (CPS / Cesu), Carlos Moreira e eu.
Optamos por um modelo diferente, convidando estudantes de vários semestres de Gestão Empresarial.
Nesse projeto, os 57 alunos brasileiros e holandeses abordaram a mesma situação problema sobre feedback (do PCI de 2018).
Durou 6 semanas, com grande enfoque na parte cultural da comunicação, da liderança e da gestão de equipes.
Agora, estou desenvolvendo um projeto com o Symbiosis Centre for Management Studies (Índia).
As expectativas são as melhores possíveis, pois percebo uma química muito grande entre os professores Nehajoan Panackal (Symbiosis), Rafaeli Begalli Danilo Sorroce (Fatec Sumaré) e eu.
Não importa se um projeto durar 4 ou 7 semanas, é importante uma preparação estratégica prévia, para entender as expectativas, interesses e ansiedades dos professores.
Ao mesmo tempo, compreender os alunos: quantos são, conhecimento do idioma, idade, o que já estudaram, semestre de curso.
Outro fator importante é a afinidade entre as disciplinas dos professores envolvidos.
No meu caso, leciono Gestão de Pessoas e Comportamento Organizacional, disciplinas que possuem aderência com várias áreas da administração.
A principal dica é não se ater somente à \textbf{ementa} da disciplina, mas ter uma visão ampla sobre as competências a serem desenvolvidas: cultura, comunicação, gestão de conflitos e trabalho em equipe.
Daí a relevância do trabalho da equipe dos PCIs / Cesu, ao encontrar parcerias personalizadas para cada professor ou professora que se interessar em participar de um PCI.
Leia a íntegra do depoimento em: https://cesu.cps.sp.gov.br/professor-da-fatec-americana-compartilha-experiencia-em-pcis/

% matched lemmas: ementa
\end{textsample}
